


\section[Spatial distribution of the kinematically cold and warm components]{Spatial distribution of the kinematically cold and warm components}
\label{sec:phase_spatial_distribution}

Figure~\ref{fig:individual_galaxies} shows the total, kinematically cold and kinematically warm \HI\ surface-density maps for the individual galaxies. Both components extend across much of the mapped atomic discs. Both components are present in clumps and ridges within the extended \HI\ emission.


The kinematically cold component contributes a median of about 45\% of the mapped \HI\ mass, while the warm component contributes the remaining 55\%. Within the optical disc ($r<r_{25}$), the warm component dominates the \HI\ budget, although the cold component remains a substantial part of the atomic gas.

The cold component has a more filamentary appearance than the warm component, with thin, elongated structures standing out against the more diffuse atomic emission. These elongated structures raise the question of whether the cold component is preferentially found in spiral arms. Such a preference could arise because compression in spiral arms can increase the gas density and promote cooling \citep{KimSpiral2008}. After tracing the spiral arms in nine spiral galaxies with visible arm structure in our sample following the methodology used by \citet{Foyle2010}, which used the 3.6 $\mu$m emission to trace stellar arms, we compared the fraction of the narrow component found in the arms and interarm regions in these galaxies. We found that the kinematically cold component is not consistently higher in the arm regions.
%To trace spiral arms in nine spiral galaxies with visible arm structure in our sample, we follow the methodology used by \citet{Foyle2010}, which used the 3.6 $\mu$m emission to trace stellar arms. We compare the fraction of the narrow component found in the arms and interarm regions in these galaxies. We found that the kinematically cold component is not consistently higher in the arm regions.

%maybe move the below to HI-H2 transition
Models of \HI\ thermal balance predict that cold, dense gas can coexist with warm, diffuse gas over a range of thermal pressures \citep{Wolfire2003}. At similar thermal pressure, the cold gas is denser because of its lower temperature. Simulations show that increasing the density of warm atomic gas can allow it to cool and condense into cold structures \citep{Saury2014}. Cold atomic gas can then supply material for molecular-cloud formation \citep{Krumholz2009Transition}.
%H$_2$ forms mainly on the surfaces of dust grains, and its formation is faster at higher gas densities \citep{GloverMacLow2007}. However, ultraviolet radiation can dissociate the molecules. Molecular gas can accumulate where dust and H$_2$ itself absorb enough of this radiation to protect the molecules. The transition from atomic to molecular gas therefore depends on the gas density, the shielding column and the dust abundance \citep{Krumholz2009Transition,Sternberg2014}.
In the simulations of \citet{GloverMacLow2007}, the initially atomic gas cools rapidly, and turbulent compression produces dense regions in which H$_2$ forms more rapidly. The connection to molecular-cloud formation therefore comes from the high densities that cold atomic gas can reach, together with sufficient shielding. This provides a reason to look for an association between the narrow \HI\ component and CO emission.

We found that both atomic components can be enhanced near CO-bright emission, but the cold fraction is not consistently higher there. The clearest enhancement occurs in NGC 5194, where the area-weighted mean cold fraction is about 45\% in the selected CO-bright regions, compared with 29\% in nearby CO-faint regions. In NGC 6946, both components increase near bright CO while the cold fraction remains about 30\%. The rest of the studied sample provides few suitable matched regions to enable a meaningful comparison between CO-bright and CO-faint emission. We therefore find no consistent evidence that narrow \HI\ is generally more concentrated near CO emission than broad \HI.

The absence of a consistent enhancement in the cold fraction does not rule out the fact that cold atomic gas is supplying material for molecular-cloud formation. In atomic-to-molecular transition models, atomic gas remains in shielding layers around molecular interiors, while additional gas can become molecular \citep{Krumholz2009Transition,Sternberg2014}. Once hydrogen becomes H$_2$, it no longer contributes to the \HI\ emission, so an increasing molecular reservoir need not be accompanied by an increasing narrow-\HI\ contribution. Furthermore, H$_2$ can survive where CO is dissociated \citep{Wolfire2010}. Our CO-faint comparison regions are therefore not necessarily free of H$_2$. These effects could weaken a spatial contrast between the CO-bright and CO-faint regions. Thus, our comparison does not establish whether the narrow \HI\ component preferentially traces regions where atomic gas becomes molecular. Locating the \HI-to-H$_2$ transition regions in our sample is difficult with our data, since a region may contain H$_2$ without detectable CO emission \citep{Wolfire2010}.
